% ====================================================================
%  CHAPTER: PROPOSED SOLUTION
% ====================================================================
\chapter{Proposed Solution}
\label{ch:proposed-solution}


Chapter~\ref{ch:fused-design-space} characterized the fused-layer
mapping design space: the proliferation of dimensions
(Section~\ref{sec:fds-dim-proliferation}), the five operand categories
(Section~\ref{sec:fds-five-operands}), the cascading tile dependencies
(Section~\ref{sec:fds-tile-deps}), and the combinatorial growth of the
map space
(Section~\ref{sec:fds-cardinality}).
This chapter presents how we \textbf{extended FactorFlow} to model,
evaluate, and navigate that space.
The extension rests on three conceptual pillars, presented in
Section~\ref{sec:ps-modeling}:
\emph{aliasing}, which captures the two-sided view of every
intermediate activation, produced by one layer, consumed by the next;
\emph{decoupling}, which guarantees each layer's tiling remains locally independent; and a
\emph{producer--consumer feasibility pruning} rule, which eliminates
infeasible mappings by enforcing the producer--consumer tile-size
inequality at every hierarchy level.

Section~\ref{sec:ps-abstraction} describes the concrete data structures
that extend FactorFlow's workload, architecture, and mapping
specifications to multi-layer fusion.
Section~\ref{sec:ps-cost-model} develops the per-layer cost-model
extensions, memory operations, energy, and latency with
bandwidth-limited stall detection, that generalize the single-layer
model of Section~\ref{sec:bg-cost-model}.
Section~\ref{sec:ps-implementation} discusses the implementation: the
overall code structure, the \texttt{moveFactor} mapper primitive, and
the current limitations.


% ====================================================================
%  §2  MODELING FUSED COMPUTATION
% ====================================================================
\section{Modeling Fused Computation}
\label{sec:ps-modeling}

In our work, a fused group of $F$ consecutive convolution layers executes as a
single, unified loop nest over a combined set of dimensions
$\mathcal{D}$.
Three modeling concepts govern how we represents the
inter-layer relationships within this unified nest.

% --------------------------------------------------------------------
\subsection{Dimension Aliasing}
\label{sec:ps-aliasing}

When layers are fused, each intermediate activation $A_i$
($0 \le i \le F\!-\!2$) serves a dual role: it is the \emph{output} of
layer~$i$ and the \emph{input} of layer~$i\!+\!1$.
The central modeling question is: which loop variables should index
this shared tensor?


\subsubsection{The Problem: Naming Intermediate Dimensions}

In a single-layer convolution the output is indexed by its own native
dimensions ($Z, P, Q$) and the input is indexed through affine
dim-sums ($[P, R], [Q, S]$), as discussed in
Section~\ref{sec:normal-affine}.
When two convolutions are fused the intermediate tensor connects them:
the output dimensions of the first convolution become the input
dimensions of the second.

Expressing the input of the second
convolution in terms of the final output would require
\emph{nested dim-sums}.
For a two-layer fusion ($F\!=\!2$) where the second convolution reads
at position $(p + r_1,\; q + s_1)$ and the first convolution must
produce tiles large enough to cover that range, the first
convolution's input would be indexed as:
%
\[
  \text{In}[\,c_0,\;
    \underbrace{(p + r_1)}_{\text{int.\ height}} + r_0,\;
    \underbrace{(q + s_1)}_{\text{int.\ width}} + s_0\,]
  \;\equiv\;
  \text{In}[\,c_0,\;
    [\,[P, R_1],\, R_0\,],\;
    [\,[Q, S_1],\, S_0\,]\,].
\]
%
For an $F$-layer fusion this nesting would grow to depth $F\!-\!1$,
producing constructs such as
$[[\cdots[[P, R_{F-1}], R_{F-2}], \cdots], R_0]$.

Nested dim-sums make it impossible to iterate over
the intermediate height and width \emph{independently} per layer: those
dimensions have no names of their own and therefore cannot appear as
standalone loops in the mapping.
This has two severe consequences:

\begin{enumerate}
  \item \textbf{No per-layer mapping freedom.}
        Without independent intermediate dimensions the mapping cannot
        express a different tiling or dataflow for each layer.
        Every loop that moves the first convolution simultaneously
        moves the second, because the only available variables ($P$,
        $Q$, $R_1$, $S_1$, $R_0$, $S_0$) are shared across the
        nested dim-sum.
        There is no way to place, say, a weight dimension of the first
        layer at one position in the dataflow while placing a weight
        dimension of the second layer at a different position ---
        the two layers cannot have independent stationarities
        (Section~\ref{sec:data-reuse}).

  \item \textbf{Forced back-to-back MAC execution.}
        Because no loop variable can advance the first convolution
        without simultaneously advancing the second, the two MAC
        operations are locked into a \emph{back-to-back} schedule:
        at every innermost iteration, $\text{MAC}_0$ produces one
        intermediate element and $\text{MAC}_1$ immediately consumes
        it.
        The mapping has no mechanism to express, for instance,
        performing several iterations of $\text{MAC}_0$ (accumulating
        a partial tile of the intermediate) before switching to
        $\text{MAC}_1$ to consume that tile --- even though such a
        schedule could improve weight reuse or reduce partial-sum
        traffic.
        More broadly, neither a layer-by-layer schedule (complete all
        $\text{MAC}_0$ iterations before any $\text{MAC}_1$) nor any
        intermediate granularity is representable; the only schedule
        is strictly element-wise interleaving.
\end{enumerate}


\subsubsection{The Solution: Concrete Intermediate Dimensions}
\label{sec:ps-alias-concrete}

Aliasing resolves this by introducing a \emph{new set of loop
variables} for each intermediate tensor.
The idea is to give the intermediate its own concrete dimension names,
for instance $Y_0$ and $X_0$ for the height and width of the first
intermediate activation, and to register them in the coupling \ref{sec:bg-coupling} so
that:
%
\begin{itemize}
  \item the \textbf{producer} (layer~$0$) indexes the intermediate
        using these names as \emph{native} dimensions, just as a
        single-layer convolution indexes its output;
  \item the \textbf{consumer} (layer~$1$) indexes the \emph{same}
        tensor through a \emph{dim-sum} involving its own variables,
        treating the intermediate, and the physical tensor, as its input.
\end{itemize}
%
Formally, for the connection between layer~$i$ and layer~$i\!+\!1$
and the tensor $A_i$, we defines two coupling entries:
%
\begin{itemize}
  \item A \textbf{producer coupling} (\texttt{int\_out}$[i]$) that
        indexes $A_i$ with native dimensions:
        \[
          A_i[\,z_i,\; y_i,\; x_i\,].
        \]
  \item A \textbf{consumer coupling} (\texttt{int\_in}$[i]$) that
        indexes the same tensor through affine dim-sums:
        \[
          A_i[\,c_{i+1},\; y_{i+1}\!+\!r_{i+1},\;
                           x_{i+1}\!+\!s_{i+1}\,].
        \]
\end{itemize}
%
The dimensions $Y_i$ and $X_i$ are now concrete: they appear in the
unified loop nest, can be tiled, reordered, and parallelized just
like any other dimension.
At the same time, the consumer sees the intermediate through its own
names ($Y_{i+1}, R_{i+1}$), so that each layer's coupling
remains a one-level list, no nesting required.


\subsubsection{Two-Layer Convolution Example}
\label{sec:ps-alias-example}

Consider the simplest non-trivial case: a fusion of two
$2$-D convolutions ($F\!=\!2$).
The complete coupling, as defined in the implementation, is:

\begin{center}
\renewcommand{\arraystretch}{1.25}
\begin{tabular}{lll}
  \toprule
  \textbf{Operand}  & \textbf{Coupling entry} & \textbf{Indexing} \\
  \midrule
  \texttt{in}
      & $[\,C_0,\; [Y_0, R_0],\; [X_0, S_0]\,]$
      & $\text{In}[\,c_0,\; y_0\!+\!r_0,\; x_0\!+\!s_0\,]$ \\[3pt]
  \texttt{w}$[0]$
      & $[\,Z_0,\; C_0,\; R_0,\; S_0\,]$
      & $W_0[\,z_0,\; c_0,\; r_0,\; s_0\,]$ \\[3pt]
  \texttt{int\_out}$[0]$
      & $[\,Z_0,\; Y_0,\; X_0\,]$
      & $A_0[\,z_0,\; y_0,\; x_0\,]$ \\[3pt]
  \texttt{int\_in}$[0]$
      & $[\,C_1,\; [P, R_1],\; [Q, S_1]\,]$
      & $A_0[\,c_1,\; p\!+\!r_1,\; q\!+\!s_1\,]$ \\[3pt]
  \texttt{w}$[1]$
      & $[\,Z_1,\; C_1,\; R_1,\; S_1\,]$
      & $W_1[\,z_1,\; c_1,\; r_1,\; s_1\,]$ \\[3pt]
  \texttt{out}
      & $[\,Z_1,\; P,\; Q\,]$
      & $\text{Out}[\,z_1,\; p,\; q\,]$ \\
  \bottomrule
\end{tabular}
\end{center}

\noindent
The twelve unified dimensions are
$\{C_0,\, Y_0,\, X_0,\, R_0,\, S_0,\, Z_0,\, C_1,\, R_1,\, S_1,\,
Z_1,\, P,\, Q\}$, and the two MAC operations described by this
coupling are:
%
\begin{align}
  \text{MAC}_0: \quad A_0[\,z_0,\, y_0,\, x_0\,]
    &\mathrel{+}= W_0[\,z_0,\, c_0,\, r_0,\, s_0\,]
       \cdot \text{In}[\,c_0,\, y_0\!+\!r_0,\, x_0\!+\!s_0\,],
  \label{eq:mac0-alias} \\[4pt]
  \text{MAC}_1: \quad \text{Out}[\,z_1,\, p,\, q\,]
    &\mathrel{+}= W_1[\,z_1,\, c_1,\, r_1,\, s_1\,]
       \cdot A_0[\,c_1,\, p\!+\!r_1,\, q\!+\!s_1\,].
  \label{eq:mac1-alias}
\end{align}

\noindent
The key property is visible in the coupling table: the rows for
\texttt{int\_out}$[0]$ and \texttt{int\_in}$[0]$ index the
\emph{same} tensor~$A_0$ under \emph{different} names.
Each pair of dimensions that addresses the same axis of an
intermediate tensor constitutes an \emph{alias}:

\begin{itemize}
  \item \textbf{Channel alias.}
        $Z_0$ (producer: number of output filters) $\leftrightarrow$
        $C_1$ (consumer: number of input channels).
        Both index the channel axis of~$A_0$.

  \item \textbf{Spatial aliases.}
        $Y_0$ (producer: output height) $\leftrightarrow$
        $[P, R_1]$ (consumer: affine dim-sum addressing input height);
        $X_0$ (producer: output width) $\leftrightarrow$
        $[Q, S_1]$ (consumer: affine dim-sum addressing input width).
\end{itemize}

\noindent
In a general $F$-layer fusion each intermediate connection~$i$
($0 \le i \le F\!-\!2$) introduces the same three aliases, with
the last connection using the global output variables $P, Q$ and the
first connection using the global input coupling.

Aliasing is \emph{implicit} in the data structures: no explicit
pointer links $Z_i$ to $C_{i+1}$ or $Y_i$ to $Y_{i+1} + R_{i+1} - 1$.
Instead, the feasibility constraints are enforced by the producer--consumer
feasibility pruning (Section~\ref{sec:ps-pruning}), which is
evaluated after every candidate mapping modification.


% --------------------------------------------------------------------
\subsection{Dimension Decoupling}
\label{sec:ps-decoupling}

Aliasing gives each intermediate tensor concrete dimension names that
appear in the unified loop nest.
A fundamental consequence of this naming is that the dimensions of
different layers become \textbf{disjoint}, leading to a property we
call \emph{decoupling}: iterating over any dimension of one layer
leaves all other layers \emph{stationary}.


\subsubsection{Per-Layer Dimension Sets}
\label{sec:ps-decouple-sets}

The unified dimension set
$\mathcal{D}$ (Equation~\ref{eq:fused-dims}) can be partitioned into
per-layer subsets by collecting, for each layer, every dimension that
appears in any operand coupling involving that layer.
For the two-layer example of
Equations~\ref{eq:mac0-alias}--\ref{eq:mac1-alias}:
%
\begin{align}
  \mathcal{D}_0 &= \underbrace{\{C_0, Y_0, R_0, X_0, S_0\}}_{\texttt{in}}
                 \;\cup\;
                 \underbrace{\{Z_0, C_0, R_0, S_0\}}_{\texttt{w}[0]}
                 \;\cup\;
                 \underbrace{\{Z_0, Y_0, X_0\}}_{\texttt{int\_out}[0]}
                 \;=\;
                 \{C_0, Y_0, X_0, R_0, S_0, Z_0\},
  \label{eq:D0-dims} \\[6pt]
  \mathcal{D}_1 &= \underbrace{\{C_1, P, R_1, Q, S_1\}}_{\texttt{int\_in}[0]}
                 \;\cup\;
                 \underbrace{\{Z_1, C_1, R_1, S_1\}}_{\texttt{w}[1]}
                 \;\cup\;
                 \underbrace{\{Z_1, P, Q\}}_{\texttt{out}}
                 \;=\;
                 \{Z_1, P, Q, R_1, S_1, C_1\}.
  \label{eq:D1-dims}
\end{align}
%
Crucially, $\mathcal{D}_0 \cap \mathcal{D}_1 = \emptyset$: the two
sets share \emph{no} dimension name.
This disjointness is precisely what aliasing achieves, $Z_0$ and
$C_1$ refer to the same physical channel axis of $A_0$, but they are
syntactically distinct loop variables, each belonging to exactly one
layer's set.

In the general $F$-layer case, an intermediate layer~$i$
($0 < i < F\!-\!1$) collects its dimensions from the weight coupling,
the preceding \texttt{int\_in}, and the succeeding \texttt{int\_out}:
\[
  \mathcal{D}_i
  \;=\;
    \underbrace{\{C_i, Y_i, R_i, X_i, S_i\}}_{\texttt{int\_in}[i\!-\!1]}
  \;\cup\;
  \underbrace{\{Z_i, C_i, R_i, S_i\}}_{\texttt{w}[i]}
  \;\cup\;
  \underbrace{\{Z_i, Y_i, X_i\}}_{\texttt{int\_out}[i]}.
\]
Each $\mathcal{D}_i$ is disjoint from every
$\mathcal{D}_{j}$ ($j \neq i$) because the subscript index supplied
by aliasing ensures that no dimension name is reused across layers.


\subsubsection{Stationarity and Independent MAC Counting}
\label{sec:ps-decouple-stationarity}

At the innermost point of the unified loop nest, all operands are
indexed down to a single element.
In a two layer fusion, two possible MAC operations exist, one per layer, among these
elements:
\[
  \text{MAC}_0:\; A_0 \mathrel{+}= W_0 \cdot \text{In},
  \qquad
  \text{MAC}_1:\; \text{Out} \mathrel{+}= W_1 \cdot A_0.
\]

\noindent
Because $\mathcal{D}_0$ and $\mathcal{D}_1$ are disjoint, advancing
any dimension in $\mathcal{D}_0$ changes the indices of
$\text{MAC}_0$'s operands while leaving every index of $\text{MAC}_1$
\emph{fixed}.
From the perspective of $\text{MAC}_1$, nothing has moved; the
second convolution is \emph{stationary}.
Symmetrically, iterating over any dimension in $\mathcal{D}_1$ leaves
$\text{MAC}_0$ stationary.

\paragraph{Example for a 2 layer fusion.}
Consider iterating over $X_0$ (the intermediate width, a native
dimension for layer~$0$).
Each step of the $X_0$ loop advances the write position
$A_0[\ldots, x_0]$ and the read position
$\text{In}[\ldots, x_0\!+\!s_0]$; new $\text{MAC}_0$ operations are
performed.
Meanwhile, $X_0 \notin \mathcal{D}_1$: neither  $A_0$, $W_1$ as read
by the consumer (which uses the dim-sum $[Q, S_1]$, not $X_0$), nor
$\text{Out}$ change their indices, so no $\text{MAC}_1$ operation is triggered.

Conversely, iterating over $Q$ (the output width, a layer-$1$
dimension) advances the consumer's read position
$A_0[\ldots, q\!+\!s_1]$ and the output write position
$\text{Out}[\ldots, q]$; new $\text{MAC}_1$ operations are performed.
Since $Q \notin \mathcal{D}_0$, the producer's indices remain fixed
and no $\text{MAC}_0$ is executed.

\medskip
\noindent
Decoupling has two consequences for the mapper:
\begin{enumerate}
  \item The dataflow, loop ordering, can be different for each layer
  \item The tiling can be different for each layer
\end{enumerate}


\subsubsection{The Need for Pruning}

Decoupling gives the mapper freedom to tile and reorder each layer's
dimensions independently, but it introduces a risk: the unified loop
nest no longer \emph{automatically} guarantees that every intermediate
value is computed before it is consumed.

Because $A_0$ is viewed from two different perspectives, written at
position $(z_0, y_0, x_0)$ by the producer and read at
$(c_1, p\!+\!r_1, q\!+\!s_1)$ by the consumer, the mapper might
choose a loop ordering or tiling that causes the consumer to request a
position not yet written by the producer.
A na\"ive non-sense mapping could, for instance, tile $Z_0$ more finely than
$C_1$, forcing the consumer to expect more channels than the producer
has completed.

This is the fundamental motivation for the producer--consumer
feasibility pruning described in the section
(Section~\ref{sec:ps-pruning}): a constraint check that verifies, at
every hierarchy level, that the producer's tile sizes are large enough
to cover the consumer's access pattern, ensuring that the data
dependencies imposed by aliasing are never violated.

\noindent








% --------------------------------------------------------------------
\subsection{Map-Space Growth under Fusion}
\label{sec:ps-mapspace}

Section~\ref{sec:bg-mapspace} derived the map-space cardinality
formula for a single-layer mapping
(Equation~\ref{eq:ms-cardinality}).
It is useful to apply this formula to the fused multi-layer configuration
with aliasing and decoupling to quantify how the design space grows with the number of fused
layers, and motivates the need of pruning and apply heuristics on the exploration space.




% ------------------------------------------------------------
\paragraph{Applying the Cardinality Formula}
\label{sec:fds-cardinality}

Recalling the map-space cardinality (Equation~\ref{eq:ms-cardinality}):
\[
|\mathcal{MS}|
\;=\;
\underbrace{(|D|!)^{L_m}}_{\text{loop permutations}}
\;\cdot\;
\underbrace{
  \prod_{d \in D}\;\prod_{p \mid d}
    \binom{v_p(d) + L - 1}{L - 1}
}_{\text{factor allocations}}.
\]

\paragraph{Dimension Proliferation with Aliasing and Decoupling}
\label{sec:ps-dim-proliferation}

Aliasing and decoupling do not merely rename existing axes: they \emph{creates new
loop variables} that the mapper must tile and order.
For each layer added to the fusion, the coupling introduces six
fresh dimensions:
%
\begin{itemize}
  \item Three \textbf{aliased dimensions} ($Y_i$, $X_i$, $Z_i$) that
        name the intermediate activation's height, width, and channel
        axes.
        Physically, these axes already exist, the intermediate tensor
        has a height, a width, and a depth that could be indexed simply with a nested dim-sum.
        However, without aliasing these axes have no independent names
        and cannot appear as standalone loops (recall the nested
        dim-sums of Section~\ref{sec:ps-alias-concrete}).
        Once aliased, they become concrete dimensions visible to the
        mapper, each admitting its own tiling factors and its own
        position in the dataflow ordering at every level.
  \item Three genuinely new dimensions ($C_i$, $R_i$, $S_i$) that
        belong exclusively to the new layer's weight coupling.
        They introduce additional degrees
        of freedom for weight tiling and loop ordering (dataflow).
\end{itemize}
%
The total dimension count therefore grows by six for every fused
layer, yielding the $|\mathcal{D}| = 6F$ formula of
Equation~\ref{eq:fused-dims}.

The map-space impact is substantial.
Each new dimension contributes additional permutations at every memory
level (widening the dataflow search) and additional factor-allocation
choices across levels (deepening the tiling search).


\paragraph{An example: Fused 11-layer configuration.}
For the 10-layer fusion workload validated in Chapter~\ref{ch:validation},
the unified loop nest has $|D_{\text{fused}}| = 66$ dimensions, 6 new dimensions introduced for every layer fused.
(Equation~\ref{eq:fused-dims}).
If, for instance, we want to execute this workload on Eyeriss-like architecture: it has $L_m = 5$ memory levels (DRAM,
Global Buffer, InputRegister, WeightRegister, IntermediateOutputRegister) and $L = 7$ total levels (5~memory + 2~spatial
fanout levels).

The permutation term becomes:
%
\begin{equation}
\label{eq:fused-perms}
(66!)^{5}.
\end{equation}
%
Even a rough estimate shows the scale:
$66! \approx 5.4 \times 10^{92}$, so
$(66!)^5 \approx 10^{460}$, compared to $720^4 \approx 10^{11}$ for
the single-layer case.
The permutation term alone grows by over \textbf{400 orders of
magnitude}.

The factor-allocation term also grows substantially.
With 66 dimensions instead of~6,
even if each dimension has only a single prime copy \ref{sec:bg-mapspace}
($v_p(d) = 1$), the per-dimension allocation count is
$\binom{1 + 7 - 1}{7 - 1} = 7$, and across 66 dimensions:
\[
7^{66} \approx 10^{55}.
\]
With larger extents (more prime copies), this term grows further.

\paragraph{Combined estimate.}
Combining both terms, the fused map-space exceeds
$10^{460} \times 10^{55} = 10^{515}$ even under conservative
assumptions.

Table~\ref{tab:mapspace-comparison} summarizes the comparison.

\begin{table}[htbp]
  \centering
  \caption{Map-space size comparison: single layer vs.\ 11-layer
           fusion on a DepFiN-like architecture.}
  \label{tab:mapspace-comparison}
  \small
  \begin{tabular}{l r r}
    \toprule
    \textbf{Quantity}
      & \textbf{Single layer}
      & \textbf{11-layer fusion} \\
    \midrule
    Dimensions ($|D|$)          & 6     & 66 \\
    Memory levels ($L_m$)       & 5& 5\\
    Total levels ($L$)          & 7& 7\\
    Permutation term            & $\sim\!10^{11}$
                                & $\sim\!10^{460}$\\
    Factor-allocation term      & $\sim\!10^{13}$--$10^{20}$
                                & $\sim10^{55}$\\
    \addlinespace
    \textbf{Total map-space}    & $\sim\!10^{24}$--$10^{32}$
                                & $\sim10^{515}$\\
    \bottomrule
  \end{tabular}
\end{table}



% --------------------------------------------------------------------
\subsection{Producer--Consumer Feasibility Pruning}
\label{sec:ps-pruning}

Aliasing (Section~\ref{sec:ps-aliasing}) establishes that intermediate
activations are shared between a producer and a consumer layer.
This sharing imposes a hard physical constraint: the consumer cannot
read a datum that the producer has not yet written.
Any mapping that violates this constraint is \emph{infeasible},  it
would cause incorrect execution regardless of the scheduling strategy
(depth-first, pipelined, or otherwise).
The following pruning rule encodes this constraint and is used to
eliminate infeasible mappings from the search space \emph{before} cost
evaluation.


\paragraph{Exactness.}
This rule is not an approximation: it encodes a necessary physical
property of any valid fused-layer schedule.
No mapping that satisfies the constraints is discarded.
The pruning is therefore \emph{exact}: it removes only infeasible
points from the map space, preserving every feasible, and potentially
optimal, mapping.


\paragraph{Per-axis decomposition.}
Each intermediate layer ~$i$ introduces three alias pairs
(Section~\ref{sec:ps-aliasing}), one for each physical axis of the
tensor~$A_i$: width $(Q_i, S_i, X_i)$, height $(P_i, R_i, Y_i)$,
and channels $(C_{i+1}, Z_i)$.
Because the three axes index orthogonal dimensions of the same tensor,
the pruning constraint can be evaluated \emph{independently} for each
triplet: satisfying the width constraint
(Equation~\ref{eq:ps-heuristic-width}) neither implies nor precludes
satisfaction of the height constraint
(Equation~\ref{eq:ps-heuristic-height}), and likewise for the channel
constraint (Equation~\ref{eq:ps-heuristic-channel}).
The implementation exploits this decomposition by calling the same
single-triplet validator specifically for the axis that was moved by move-factor with respect to the starting valid configuration.

\subsubsection{Cross-Level Validation}
\label{sec:ps-cross-level}

The core of the pruning logic is the observation that, in a
multi-memory level loop nest, the producer dimension (e.g.: $X$) may be tiled
across several hierarchy levels, while the consumer dimensions (e.g.: $Q$
and $S$) that fall \emph{between} two consecutive $X$ loop instances
can also belong to different levels.
The constraint must therefore be checked not at each memory level in
isolation, but at each boundary defined by the $X$ loops in the
nest, a boundary that can span multiple memory levels.

\paragraph{Zone definition.}
Consider a loop order that contains $K$ instances of the producer
dimension $X$: $x_1, x_2, \dots, x_K$ (from outermost to innermost).
Each consecutive pair $(x_k, x_{k+1})$ defines a \emph{zone} — the
set of all loops strictly between $x_k$ and $x_{k+1}$ in the nest.
Within a zone, the consumer loops $Q$ and $S$ may appear at any
memory level; the zone itself typically spans the boundary between
two or more hierarchy levels.

\paragraph{Accumulated iterations.}
For each zone, the validator computes:
\begin{itemize}
  \item $Q_{\text{zone}}$: the product of iterations of all $Q$ loops
        in the zone (regardless of which memory level they belong to).
  \item $S_{\text{zone}}$: the product of iterations of all $S$ loops
        in the zone.
  \item $X_{\text{cum}}$: the product of iterations of the current $X$
        loop and all $X$ loops inner to it (i.e.\ the cumulative
        producer tile size from this boundary inward).
\end{itemize}
The constraint then requires:
\begin{equation}
\label{eq:ps-zone-check}
  Q_{\text{zone}} + S_{\text{zone}} - 1 \;\le\; X_{\text{cum}}.
\end{equation}

\noindent
At the \emph{outermost} $X$ loop, the zone extends to the top of the
nest, so $Q_{\text{zone}}$ and $S_{\text{zone}}$ become the
\emph{total global} iterations of $Q$ and $S$ across all levels, and
$X_{\text{cum}}$ becomes the total extent of $X$.
This outermost check ensures that the full-dimension constraint
$Q + S - 1 \le X$ is satisfied globally.


\subsubsection{Initial Condition}

An additional constraint addresses the very bottom of the loop nest.
The innermost loop with more than one iteration must be a
\emph{producer} (input spatial) dimension, not a consumer or weight
dimension.
This guarantees that the first non-trivial iterations of the nest
produce intermediate data before any consumer iteration attempts to
read it.
If a consumer dimension $Q$ or $S$ has non-trivial iterations below
the innermost $X$, the mapping is rejected immediately.



\subsubsection{Integration with the Mapper}

Every candidate mapping move (Section~\ref{sec:ps-movefactor}) can
trigger a call to the pruning validator.
If the move results in a mapping that violates any of
Equations~\ref{eq:ps-heuristic-width}--\ref{eq:ps-heuristic-channel},
the move is rolled back and marked as infeasible.
Because the pruning eliminates only infeasible mappings, the mapper
is guaranteed to retain every valid candidate.


% ====================================================================
%  §3  EXTENDING THE FACTORFLOW ABSTRACTION
% ====================================================================
\section{Extending the FactorFlow Abstraction}
\label{sec:ps-abstraction}

FactorFlow describes a mapping problem through three inputs: a
\emph{workload}, an \emph{architecture}, and a \emph{mapping}.
Each must be extended to support multi-layer fusion.
The extensions are designed to be \emph{backward-compatible}: a
single-layer problem is a special case with $F = 1$ and no
intermediate operands.

% --------------------------------------------------------------------
\subsection{Workload Specification}
\label{sec:ps-workload-spec}

A workload is defined by two objects: a \textbf{Shape} \ref{sec:bg-shape} and a
\textbf{Coupling} \ref{sec:bg-coupling}.

The \texttt{Shape} is a dictionary mapping every dimension name to its
size.
For a fused group of $F$ convolution layers operating on 2-D feature
maps, the shape contains entries for:
\begin{itemize}
  \item $C_0$: input channels; and, for the first layer, $Y_0$ and $X_0$
        (the intermediate spatial extent produced by layer~0).
  \item For each layer $i = 0, \dots, F\!-\!1$: $Z_i$, $C_i$, $R_i$, $S_i$
        (output channels, input channels, filter height, filter width).
  \item For each intermediate $i = 1, \dots, F\!-\!2$: $Y_i$ and $X_i$
        (spatial height and width of the intermediate activation produced
        by layer~$i$).
  \item $P$, $Q$: final output height and width; $Z_{F-1}$: final output
        channels.
\end{itemize}

\noindent
The \texttt{Coupling} encodes which dimensions index which operand (or tensor).
Its constructor accepts \emph{five} coupling lists, one for each
operand category introduced in Section~\ref{sec:fds-five-operands}:

\begin{center}
\small
\begin{tabular}{lll}
  \toprule
  \textbf{Argument} & \textbf{Type} & \textbf{Semantics} \\
  \midrule
  \texttt{in\_coupling}
      & \texttt{list}
      & Dims indexing the global input \\
  \texttt{w\_coupling}
      & \texttt{dict[\,int $\to$ list\,]}
      & Per-layer weight dims \\
  \texttt{out\_coupling}
      & \texttt{list}
      & Dims indexing the global output \\
  \texttt{int\_out\_coupling}
      & \texttt{dict[\,int $\to$ list\,]}
      & Per-connection producer dims \\
  \texttt{int\_in\_coupling}
      & \texttt{dict[\,int $\to$ list\,]}
      & Per-connection consumer dims \\
  \bottomrule
\end{tabular}
\end{center}

\noindent
The dictionary key for \texttt{w\_coupling} is the layer index
$i \in \{0, \dots, F\!-\!1\}$; for the intermediate couplings it is
the connection index $i \in \{0, \dots, F\!-\!2\}$.
Internally, the \texttt{Coupling} class stores both the raw (nested)
and a flattened version of each coupling list.
The flattened form is used for efficient look-ups during MOPs
computation and constraint checking; the nested form preserves the
affine structure for tile-size calculations.

\paragraph{Affine indexing in couplings.}
A nested list inside a coupling entry denotes affine (summed) indexing.
For example, $[\texttt{Y1},\; \texttt{R1}]$ in
$\texttt{int\_in\_coupling[0]}$ indicates activation indexing of the
form $y_1 + r_1$, matching the sliding-window access pattern of a
convolution.
Optional stride dictionaries (\texttt{in\_strides},
\texttt{w\_strides}, etc.) allow each contributing dimension to carry a
multiplicative coefficient, generalising the model to strided
convolutions (cf.\ Section~\ref{sec:normal-affine}).

\paragraph{Example: two-layer convolution.}
The coupling for a two-layer fusion ($F = 2$) provides a concrete
illustration (cf.\ Section~\ref{sec:multi-layer-workload}):

\begin{center}
\small
\begin{tabular}{@{}l@{\;:\;\;}l@{}}
  \texttt{dims}
    & $[\,C_0,\, Y_0,\, X_0,\, R_0,\, S_0,\, Z_0,\,
       C_1,\, R_1,\, S_1,\, Z_1,\, P,\, Q\,]$ \\[3pt]
  \texttt{in\_coupling}
    & $[\,C_0,\; [Y_0, R_0],\; [X_0, S_0]\,]$ \\
  \texttt{w\_coupling}
    & $\{0\!: [\,Z_0, C_0, R_0, S_0\,],
       \;\; 1\!: [\,Z_1, C_1, R_1, S_1\,]\}$ \\
  \texttt{int\_out\_coupling}
    & $\{0\!: [\,Z_0,\, Y_0,\, X_0\,]\}$ \\
  \texttt{int\_in\_coupling}
    & $\{0\!: [\,C_1,\; [P, R_1],\; [Q, S_1]\,]\}$ \\
  \texttt{out\_coupling}
    & $[\,Z_1,\, P,\, Q\,]$ \\
\end{tabular}
\end{center}

\noindent
Twelve dimensions, five operand categories, and the aliasing between
$\texttt{int\_out}[0]$ and $\texttt{int\_in}[0]$ is immediately
visible: $Z_0 \leftrightarrow C_1$ (channel aliasing),
$(Y_0, X_0) \leftrightarrow ([P, R_1],\, [Q, S_1])$ (spatial aliasing).

The pattern generalizes mechanically: adding a layer appends six new
dimensions ($Z_i, C_i, R_i, S_i, Y_i, X_i$ for intermediate layers,
or $P, Q$ for the last), one new \texttt{w\_coupling} entry, and one
new pair of intermediate coupling entries.


% --------------------------------------------------------------------
\subsection{Architecture Specification}
\label{sec:ps-arch-spec}

An architecture in FactorFlow is an ordered list of \emph{levels},
from outermost (e.g.\ DRAM) to innermost (Compute).
Three level types are supported:
\begin{itemize}
  \item \textbf{MemLevel} — a storage level characterized by capacity
        (bytes), per-access energy (picojoules per byte), separate read
        and write bandwidths (bytes per cycle), a bypass list, and
        optional factors and dataflow constraints.
  \item \textbf{FanoutLevel} — a spatial unrolling level with a fixed
        mesh size and the dimensions across which data is distributed.
  \item \textbf{ComputeLevel} — the innermost functional unit
        performing the MAC operation, with a fixed compute energy and
        latency (cycles per operation).
\end{itemize}

\paragraph{Five-category bypasses.}
The single-layer FactorFlow uses three bypass tags
(\texttt{in}, \texttt{w}, \texttt{out}).
The multi-layer extension adds \texttt{int\_in} and \texttt{int\_out},
enabling per-level control over which intermediate operands are stored.
A level's bypass list specifies which operand categories are
\emph{not} stored at that level — equivalently, which operands "pass
through" without incurring local access energy.

As an example, the DepFiN architecture
(cf.\ Table~\ref{tab:bypass-depfin}) specifies:
\begin{center}
\small
\begin{tabular}{ll}
  \toprule
  \textbf{Level} & \textbf{Bypassed operands} \\
  \midrule
  DRAM           & \texttt{int\_in}, \texttt{int\_out} \\
  FeatureMemory  & \texttt{w} \\
  WeightMemory   & \texttt{in}, \texttt{int\_in},
                   \texttt{int\_out}, \texttt{out} \\
  AccReg         & \texttt{in}, \texttt{w}, \texttt{int\_in} \\
  \bottomrule
\end{tabular}
\end{center}

\noindent
Critically, both \texttt{int\_in} and \texttt{int\_out} bypass DRAM:
intermediate activations never leave the chip.
They are stored exclusively in the FeatureMemory (FMEM), achieving the
inter-layer data reuse that motivates fusion
(Section~\ref{sec:fds-immediate-reuse}).

\paragraph{Dataflow and factors constraints.}
Each \texttt{MemLevel} can specify:
\begin{itemize}
  \item A \texttt{dataflow\_constraints} list, fixing the loop
        ordering at that level.  For instance, the WeightMemory in
        DepFiN constrains all weight dimensions
        ($Z_i, C_i, R_i, S_i$ for every~$i$) to be the innermost
        loops, ensuring maximal weight reuse.
  \item A \texttt{factors\_constraints} dictionary, fixing the number
        of iterations of specified dimensions.  For example,
        $\{C_i\!: 32\}$ forces the innermost weight loop on the
        channel dimension to tile by~32, while $\{R_i\!: 3,\; S_i\!: 3\}$
        places the full $3 \times 3$ filter at that level.
\end{itemize}

\noindent
The method \texttt{enforceFactorsConstraints} distributes the
constraint factors automatically upon initialization, moving the
required prime factors from the innermost level (where all factors
start) to the constrained level.


% --------------------------------------------------------------------
\subsection{Mapping Specification}
\label{sec:ps-mapping-spec}

A mapping in FactorFlow is defined by two pieces of information at
every level of the architecture:
\begin{enumerate}
  \item The \textbf{factors allocation}: the prime-factorized tiling
        factors assigned to each dimension at that level, determining
        how many iterations of each dimension execute there
        (cf.\ Section~\ref{sec:tiling}).
  \item The \textbf{dataflow}: the ordering of dimensions at that
        level, determining the loop nesting and hence which operand
        enjoys innermost reuse
        (cf.\ Section~\ref{sec:loop-ordering}).
\end{enumerate}

\noindent
The product of a dimension's factors across all levels must equal the
dimension's full size (completeness), and the tile sizes at each level
must respect the capacity and equality constraints attached to that
level (feasibility).

For fused mappings, the factors allocation spans all $|\mathcal{D}|$
dimensions — both aliased and decoupled — and the
producer--consumer pruning of Section~\ref{sec:ps-pruning} adds an
additional feasibility filter on top of the standard capacity
constraints.

\paragraph{Per-layer stationarity within a single loop nest.}
A key consequence of aliasing and decoupling
(Sections~\ref{sec:ps-aliasing}--\ref{sec:ps-decoupling}) is that the
unified dataflow implicitly encodes a \emph{different mapping}, and
hence a different \emph{stationarity}, for each layer, all within a
single loop nest.

Recall that stationarity is determined by which dimensions iterate at
a given level: an operand is stationary when the innermost loops at
that level are orthogonal to it
(Section~\ref{sec:data-reuse}).
Because decoupling makes the per-layer dimension sets disjoint
($\mathcal{D}_i \cap \mathcal{D}_j = \emptyset$,
Section~\ref{sec:ps-decouple-sets}), the dataflow ordering at any
level is \emph{filtered per layer}: at each level only the dimensions
in~$\mathcal{D}_i$ are visible to layer~$i$.
Within the unified ordering, each layer therefore ``sees'' its own
sub-sequence of loops, and the innermost element of that sub-sequence
determines that layer's stationarity at that level.

Consider, for instance, a two-layer fusion where the dataflow at a
memory level orders the dimensions as
$[\ldots,\, C_1,\, X_0,\, \ldots]$ (from outermost to innermost).
Layer~$0$ sees $X_0$ as the innermost relevant loop: the weights
$W_0[\,z_0, c_0, r_0, s_0\,]$ do not depend on $X_0$, so $W_0$ is
stationary.
Layer~$1$ sees $C_1$ as its innermost relevant loop: the output
$\text{Out}[\,z_1, p, q\,]$ does not depend on $C_1$, so the output
is stationary.
A different ordering, e.g.\
$[\ldots,\, Z_1,\, X_0,\, \ldots]$, would make
$W_1$ vary on the innermost layer-$1$ loop ($Z_1$ indexes $W_1$),
yielding an \emph{input-stationary} dataflow for layer~$1$ while
layer~$0$ remains output-stationary.
The same unified nest thus accommodates layer~$0$ in output-stationary
mode and layer~$1$ in input-stationary mode, no separate per-layer
scheduling is needed.

This property follows directly from the modeling choices described
earlier: aliasing gives each intermediate tensor its own set of
concrete loop variables (Section~\ref{sec:ps-alias-concrete}), and
decoupling guarantees that these per-layer variables are disjoint
(Section~\ref{sec:ps-decouple-sets}).
The unified dataflow is therefore strictly more expressive than $F$
independent single-layer mappings laid out in sequence: it can
reproduce any combination of per-layer stationarities while also
capturing the inter-layer data dependencies through the
producer--consumer pruning constraint.

\paragraph{Initialization sequence.}
Mapping construction follows a deterministic five-step sequence:
\begin{enumerate}
  \item \texttt{initFactors}: all prime factors of every dimension
        are placed on the innermost level.
  \item \texttt{enforceFactorsConstraints}: factors are moved outward
        to satisfy each level's equality ($=$) and lower-bound
        ($\ge$) factors constraints, using \texttt{moveFactor}
        (Section~\ref{sec:ps-movefactor}) with constraint checking
        temporarily disabled.
  \item \texttt{setupBypasses}: for each bypassed operand, the last
        non-bypassing \texttt{MemLevel} receives a pointer to the
        chain of bypassed levels it must serve, enabling its
        \texttt{MOPs()} method to account for the extra traffic.
  \item \texttt{setupSpatialLevelPointers}: each \texttt{MemLevel}
        is linked to the spatial levels (fanouts and compute) that
        follow it.
  \item \texttt{checkFactorsConstraints}: a final global validation
        ensures the entire architecture satisfies all constraints.
\end{enumerate}

\noindent
After initialization, the resulting mapping, iff feasible, serves as the starting
point for the map-space exploration engine, which navigates the factor
space using \texttt{moveFactor} (Section~\ref{sec:ps-movefactor}).

% ====================================================================
%  §4  COST-MODEL EXTENSIONS
% ====================================================================
\section{Cost-Model Extensions}
\label{sec:ps-cost-model}

The single-layer cost model (Section~\ref{sec:bg-cost-model}) computes
three quantities: memory operations (MOPs), energy, and latency.
Each must be generalized to handle $F > 1$ layers with five operand
categories.

% --------------------------------------------------------------------
\subsection{Per-Layer Memory Operations}
\label{sec:ps-mops}

The entry point \texttt{updateStats} traverses the architecture
\emph{top-to-bottom} (from DRAM toward Compute).
At each \texttt{MemLevel}, the method \texttt{MOPs()} returns the
\emph{per-instance} read and write counts for every operand category:

\begin{center}
\small
\begin{tabular}{ll}
  \toprule
  \textbf{Return value} & \textbf{Description} \\
  \midrule
  \texttt{in\_reads}
      & Reads of the global input \\
  \texttt{per\_layer\_w\_reads}$[i]$
      & Reads of weight tensor $W_i$ \\
  \texttt{per\_layer\_int\_in\_reads}$[i]$
      & Reads of $A_i$ as consumer input \\
  \texttt{per\_layer\_int\_out\_reads}$[i]$
      & Partial-sum reads of $A_i$ (accumulation) \\
  \texttt{per\_layer\_int\_out\_writes}$[i]$
      & Partial-sum writes of $A_i$ \\
  \texttt{out\_reads}
      & Partial-sum reads of the global output \\
  \texttt{out\_writes}
      & Writes of the global output \\
  \bottomrule
\end{tabular}
\end{center}

\noindent
These per-instance counts are then \textbf{scaled} to total counts by
multiplying each layer's contributions by the product of temporal and
spatial iterations relevant to that layer:
%
\begin{equation}
\label{eq:ps-scaling}
  \sigma_l^{(i)} \;=\;
    \tau_l^{(i)} \;\cdot\; s_l^{(i)},
\end{equation}
%
\noindent
where $\tau_l^{(i)} = \prod_{k > l}\,\prod_{d \in \mathcal{D}_i}
f_{k,d}$ is the product of temporal iterations of all dimensions
relevant to layer~$i$ at levels above~$l$, and $s_l^{(i)}$ is the
analogous spatial iteration count accumulated from fanout levels.
The set $\mathcal{D}_i$ contains only the dimensions coupled to
layer~$i$'s operands, for example, layer~0 is influenced only by
$\{C_0,\, Y_0,\, X_0,\, R_0,\, S_0,\, Z_0\}$.

\paragraph{Per-layer dimension filtering.}
At each level, a \texttt{dataflow\_per\_layer} dictionary is
constructed, filtering the level's full dataflow to retain only the
dimensions relevant to each layer.
This filtering ensures that a dimension belonging solely to layer~$j$
does not inflate the iteration count attributed to layer~$i$.
The method \texttt{fullLayerProduct} computes the product of factor
iterations for a specific layer at a specific level using exactly this
filtered dimension set.

\paragraph{Bypass propagation.}
When an operand bypasses a level, its reads and writes at that level
are zero.
The level \emph{above} the bypass chain is responsible for supplying
data to the level \emph{below}; FactorFlow models this by carrying
"last reads" accumulators that propagate the most recent non-bypassed
read count downward.
For intermediate operands, this mechanism applies per-layer: each
$\texttt{per\_layer\_int\_out\_reads}[i]$ and
$\texttt{per\_layer\_int\_in\_reads}[i]$ is propagated independently
through the bypass chain.

\paragraph{Drain handling.}
Writes from a lower level to a higher level (drains) are modeled
symmetrically: when a level is not bypassed for an operand, the reads
from the level above become writes into this level.
The \texttt{FREE\_DRAINS} setting controls whether drains contribute
to energy and latency; when enabled, drain reads are assumed to be
zero-cost.

% --------------------------------------------------------------------
\subsection{Energy Model}
\label{sec:ps-energy}

The energy weighted MOPs (WMOPs) metric (cf.\ Equation~\ref{eq:energy}) is
extended to a per-layer decomposition:
%
\begin{equation}
\label{eq:ps-wmops-per-layer}
  \text{WMOPs}_i \;=\;
    \sum_{l \in \text{levels}}
      e_l \cdot \bigl(R_l^{(i)} + W_l^{(i)}\bigr),
\end{equation}
%
\noindent
where $R_l^{(i)}$ and $W_l^{(i)}$ are the total reads and writes at
level~$l$ attributable to layer~$i$, and $e_l$ is the per-access
energy of level~$l$.
For a boundary layer ($i = 0$ or $i = F\!-\!1$), $R_l^{(i)}$ includes
the global input or output contributions; for intermediate layers, it
includes only the corresponding intermediate and weight operands.

The total energy is:
%
\begin{equation}
\label{eq:ps-wmops-total}
  \text{WMOPs} \;=\;
    \sum_{i=0}^{F-1} \text{WMOPs}_i
    \;+\;
    \text{WMOPs}_{\text{compute}},
\end{equation}
%
\noindent
where $\text{WMOPs}_{\text{compute}}$ accounts for the MAC energy
across all layers.
The per-layer decomposition is valuable for identifying which
layer dominates energy consumption and, in future work, for guiding
per-layer mapping optimization.


% --------------------------------------------------------------------
\subsection{Latency Model}
\label{sec:ps-latency}

Latency is computed \emph{bottom-to-top}, starting from the Compute
level and propagating upward through the memory hierarchy.
At the Compute level, each tile executes in a fixed number of cycles
(typically~1 for a single MAC).
At each \texttt{MemLevel} above, the model determines whether the
level is \textbf{compute-bound} or \textbf{memory-bound} on a
per-layer basis.

\paragraph{Per-layer ideal bandwidth.}
For each layer~$i$ at level~$l$, the ideal read bandwidth is:
%
\begin{equation}
\label{eq:ps-ideal-bw}
  b_{l,\text{rd}}^{(i)} \;=\;
    \frac{R_l^{(i)} + D_l^{(i)}}
         {\sigma_l^{(i)} \cdot T_l^{(i)}},
\end{equation}
%
\noindent
where $D_l^{(i)}$ is the drain traffic, $\sigma_l^{(i)}$ is the
scaling factor from Equation~\ref{eq:ps-scaling}, and $T_l^{(i)}$ is
the number of compute cycles per tile for layer~$i$ at level~$l$
(accounting for any stalls introduced by lower levels).
The write bandwidth $b_{l,\text{wr}}^{(i)}$ is defined analogously
using fills and updates.

\paragraph{Stall detection.}
If $b_{l,\text{rd}}^{(i)} > B_{l,\text{rd}}$ (the level's physical
read bandwidth in bytes per cycle), the level is memory-bound for
layer~$i$'s reads: the actual read-plus-drain latency becomes
%
\begin{equation}
\label{eq:ps-stall-lat}
  \Lambda_{l,\text{rd}}^{(i)} \;=\;
    \frac{R_l^{(i)} + D_l^{(i)}}
         {\sigma_l^{(i)} \cdot B_{l,\text{rd}}}.
\end{equation}
%
\noindent
Otherwise the level is compute-bound and
$\Lambda_{l,\text{rd}}^{(i)} = T_l^{(i)}$.
The same logic applies to the write side using fills and updates.
The per-layer per-level latency is:
%
\begin{equation}
\label{eq:ps-lat-per-layer}
  \Lambda_l^{(i)} \;=\;
    \max\!\bigl(
      \Lambda_{l,\text{rd}}^{(i)},\;
      \Lambda_{l,\text{wr}}^{(i)}
    \bigr),
\end{equation}
%
\noindent
and the stall cycles contributed by level~$l$ for layer~$i$ are
$\Lambda_l^{(i)} - T_l^{(i)}$.

\paragraph{Sequential layer execution.}
Under the \texttt{SEQUENTIAL\_LAYER\_EXECUTION} model the
per-level latency is the \emph{sum} of per-layer latencies:
%
\begin{equation}
\label{eq:ps-sequential-lat}
  \Lambda_l \;=\;
    \sum_{i=0}^{F-1} \Lambda_l^{(i)}.
\end{equation}
%
\noindent
This models the fact that the PE array processes one layer at a time
in a time-multiplexed fashion.
The total latency is then the maximum across all levels:
%
\begin{equation}
\label{eq:ps-total-lat}
  \Lambda \;=\;
    \max_{l \in \text{levels}} \;\Lambda_l.
\end{equation}

\paragraph{Latency for single buffered memory levels.}
A level is single buffered when it cannot performs both reads and writes at
the same time, in this situation, the latency of the level is not anymore the
maximum of latency due to the reads and to the writes, but the sum.
For instance, the FeatureMemory (FMEM) level in the DepFiN architecture is
a single buffered memory level: because it is the on-chip buffer that stores
all activations (input, intermediates, output), its read and write
latencies are \emph{summed} rather than taking the maximum:
%
\begin{equation}
\label{eq:ps-fmem-lat}
  \Lambda_{\text{FMEM}}^{(i)} \;=\;
    \Lambda_{\text{FMEM},\text{rd}}^{(i)}
    + \Lambda_{\text{FMEM},\text{wr}}^{(i)}.
\end{equation}
%
\noindent
This reflects the observation that, for the access patterns of fused
execution on DepFiN, the read and write traffic to FMEM cannot overlap
fully.
The resulting per-layer latency is then
$\max(\Lambda_{\text{FMEM}}^{(i)},\; T_l^{(i)})$, preserving the
guarantee that latency is never below the ideal compute time.

\paragraph{Compound metrics.}
The EDP (Energy-Delay Product) and Wart (weighted area-time) metrics
defined in Section~\ref{sec:compound-metrics} apply unchanged; they
simply use the extended WMOPs and latency values computed above.


% ====================================================================
%  §5  IMPLEMENTATION
% ====================================================================
\section{Implementation}
\label{sec:ps-implementation}

This section provides an overview of the FactorFlow codebase, focusing
on the modules most affected by the multi-layer extension.

% --------------------------------------------------------------------
\subsection{Code Structure}
\label{sec:ps-code-structure}

The implementation consists of approximately 5\,500~lines of Python
spread across ten principal modules:

\begin{center}
\small
\begin{tabular}{llp{7.2cm}}
  \toprule
  \textbf{Module} & \textbf{Lines} & \textbf{Responsibility} \\
  \midrule
  \texttt{factors.py}
      & 742
      & \texttt{Coupling}, \texttt{Factors} (prime-factor
        dictionaries), \texttt{Shape} classes \\
  \texttt{levels.py}
      & 2192
      & \texttt{MemLevel} (incl.\ $\sim$700-line
        \texttt{MOPs()} method), \texttt{FanoutLevel},
        \texttt{ComputeLevel} \\
  \texttt{arch.py}
      & 857
      & \texttt{Arch}: \texttt{moveFactor},
        \texttt{validate\_mapping\_heuristic},
        \texttt{enforceFactorsConstraints} \\
  \texttt{cost\_model.py}
      & 707
      & \texttt{updateStats}, \texttt{Wart}, \texttt{EDP},
        \texttt{Energy}, \texttt{Latency} \\
  \texttt{engine.py}
      & 115
      & Mapper entry point, multithreading orchestration \\
  \texttt{settings.py}
      & 187
      & Global configuration flags
        (\texttt{SEQUENTIAL\_LAYER\_EXECUTION},
        \texttt{FULLY\_CACHED}, etc.) \\
  \texttt{pruning\_check.py}& 289
      & Standalone pruning validator for unit testing \\
  \texttt{main\_cli.py}
      & 289
      & Command-line interface for specifying architecture
        and workload files \\
  \texttt{prints.py}, \texttt{utils.py}
      & ---
      & Output formatting and utility functions \\
  \bottomrule
\end{tabular}
\end{center}

\noindent
The most substantial changes for multi-layer support are concentrated
in \texttt{factors.py} (five-category \texttt{Coupling}),
\texttt{levels.py} (per-layer MOPs decomposition),
\texttt{arch.py} (feasibility pruning, per-layer constraint
management), and \texttt{cost\_model.py} (per-layer energy and
latency).

Pre-defined coupling instances in \texttt{computations.py} allow
workloads up to 10-layer fusions to be constructed by importing the
corresponding coupling object and pairing it with a \texttt{Shape}
specifying the dimension sizes.
Larger fusions can be created programmatically by extending the same
pattern.


% --------------------------------------------------------------------
\subsection{The \texttt{moveFactor} Primitive}
\label{sec:ps-movefactor}

The map-space exploration in FactorFlow is built on a single
primitive:

\begin{quote}
\texttt{moveFactor(src, dst, dim, factor, amount)}
\end{quote}

\noindent
This method transfers $\texttt{factor}^{\texttt{amount}}$ iterations
of dimension \texttt{dim} from level \texttt{src} to level
\texttt{dst}.
The operation proceeds in six steps:

\begin{enumerate}
  \item \textbf{Remove} the factor from the source level's
        prime-factor dictionary.
        If the source does not hold enough factors, the move fails
        immediately (returns \texttt{False}).
  \item \textbf{Check source constraints.}
        If the removal leaves the source level in violation of a
        capacity or equality constraint, the factor is restored and
        the move fails.
  \item \textbf{Add} the factor to the destination level.
  \item \textbf{Update tile sizes.}
        For every level between source and destination, the tile size
        for \texttt{dim} is multiplied (or divided) by
        $\texttt{factor}^{\texttt{amount}}$, reflecting the changed
        iteration count.
  \item \textbf{Check constraints on all affected levels.}
        If any level between source and destination (inclusive)
        violates its constraints after the tile-size update, the
        entire operation is rolled back: factor restored to source,
        removed from destination, tile sizes reverted.
  \item \textbf{Feasibility pruning}
        Basing on the axis of the dimension that is changed,  
        the corresponding full multi-layer producer--consumer pruning
        (Section~\ref{sec:ps-heuristic}) is evaluated on the new
        resulting mapping.
        If the axis producer--consumer pruning constraints fail for any layer, the move is fully rolled back.
\end{enumerate}

\noindent
The rollback guarantee ensures that every call to \texttt{moveFactor}
is \emph{atomic}: the architecture is either updated to a new valid
state or left entirely unchanged.
This property makes it safe to use \texttt{moveFactor} as the
exploration step in hill-climbing, breadth-first, or any other
local-search mapper without fear of accumulating partial state.

The mapper iterates through dimensions and levels, attempting
\texttt{moveFactor} calls and evaluating the cost model after each
successful move.
Different mapper variants
(\texttt{breadthfirst\_local}, \texttt{linear}, \texttt{quadratic},
\texttt{exponential}, \texttt{hybrid}) differ in how they select
which moves to try and how deeply they explore, but all rely on
\texttt{moveFactor} as their sole mutation operator.


% --------------------------------------------------------------------
\subsection{Limitations}
\label{sec:ps-limitations}

The current implementation has several limitations that constrain its
applicability:

\begin{enumerate}
  \item \textbf{No automated mapper for fused workloads.}
        The map-space exploration engine (\texttt{engine.py})
        currently supports automated search only for single-layer
        mappings.
        For multi-layer fused workloads, is performed the evaluation of the single mapping described through the architecture definition file.
        Extending the mapper to automatically explore fused mappings,
        while respecting the pruning constraints, is the
        primary direction for future work.

  \item \textbf{No inter-layer parallelism.}
        The \texttt{SEQUENTIAL\_LAYER\_EXECUTION} flag is always
        active: all layers share the same PE array and execute one at
        a time.
        Architectures that overlap computation of different layers
        (e.g.\ pipeline parallelism) are not tested.
\end{enumerate}